\section{The virtual contribution in a common function basis}
\label{sec:comp-virtual-energy-basis}
\subsection{A constant alphabet}
Use iterated integrals based at $X=0$, with the following kernels:
\begin{align}
 G(a,w;X)&=\int_0^X\omega_a(t)G(w;t),\qquad G(\varnothing;X)=1,\\
 \omega_a&=\frac{\dd X}{X-a}\quad(a=0,1,-1,-2,-4),\qquad
 \omega_r=\frac{\dd X}{\sqrt{X(X+4)}}.
\end{align}
The root is positive for $X>0$. Pure-zero words have the normalization
$G(0,\ldots,0;X)=\ln^nX/n!$. The letter $1$ occurs in combinations
whose apparent pole is canceled by an inner function vanishing at
$X=1$, as in $G(1,0;X)$ with kernel $\ln X/(X-1)$.

Define a new variable $q$, unrelated to electric charge, by
\begin{equation}
 q=\frac{\sqrt{X+4}-\sqrt X}{\sqrt{X+4}+\sqrt X},\qquad
 X=\frac{(1-q)^2}{q},\qquad 0<q<1.
\end{equation}
The physical square roots are positive.  Introduce
\begin{equation}
 \rho_\pm=\frac{1\pm\ii\sqrt3}{2},\qquad
 \gamma_\pm=\frac{3\pm\sqrt5}{2},\qquad
 \nu_a=\frac{\mathrm{d}q}{q-a}.
\end{equation}
Direct differentiation gives the complete pullback of the kernels:
\begin{align}
 \frac{\mathrm{d}X}{X}&=2\nu_1-\nu_0,\\
 \frac{\mathrm{d}X}{X+1}&=\nu_{\rho_+}+\nu_{\rho_-}-\nu_0,\\
 \frac{\mathrm{d}X}{X+2}&=\nu_{\ii}+\nu_{-\ii}-\nu_0,\\
 \frac{\mathrm{d}X}{X+4}&=2\nu_{-1}-\nu_0,\\
 \frac{\mathrm{d}X}{X-1}&=\nu_{\gamma_+}+\nu_{\gamma_-}-\nu_0,\\
 \frac{\mathrm{d}X}{\sqrt{X(X+4)}}&=-\nu_0.
 \label{eq:compton-equivalence-pullback}
\end{align}
Replacing each letter by its linear combination in these equations
converts every independently integrated word to constant-letter GPLs based at $q=1$.
The basepoint is approached from the left.  Its logarithmic convention
is fixed by the exact relation
\begin{equation}
 \ln X=2\ln(1-q)-\ln q.
\end{equation}
Writing $q=1-\delta$ gives $X=\delta^2/(1-\delta)$ and
$\ln X=2\ln\delta+O(\delta)$.  The leading coefficient of
$X$ in this relation is one, so it produces no finite logarithmic
constant.  This fixes words ending in the singular letter without an
undetermined additive constant.  All other boundary constants follow
by induction from the vanishing lower-weight words.  Complex letters
occur in conjugate pairs, with the same real path and conjugate boundary
values.  The resulting alphabet is
\begin{equation}
 \mathcal A=\{0,1,-1,\ii,-\ii,\rho_+,\rho_-,\gamma_+,\gamma_-\}.
\end{equation}
The transformed words have finite rational coefficient dictionaries.
All six kernel identities
in \eqref{eq:compton-equivalence-pullback} are checked symbolically.

\subsection{The vertex and self-energy contribution}
The entire independently integrated vertex and self-energy contribution
can already be reduced to constant letters.  In the direct channel, the
physical continuation gives
\begin{equation}
 \operatorname{Re}\operatorname{Li}_2(1+X+\ii0)
 =\frac{\pi^2}{6}-G(-1,0;X),\qquad
 \operatorname{Re}\ln(-X-\ii0)=\ln X.
\end{equation}
The imaginary part of the dilogarithm is fixed by its indicated boundary,
not by a fit to the final answer.  The interference coefficients in this
channel are real, so these real parts suffice for this contribution.
In the crossed channel the two angular endpoints are $y=X$ and
$y=X/(1+X)$.  At the latter endpoint, the substitution $y=t/(1+t)$ gives
\begin{equation}
 \frac{\mathrm{d}y}{y}=\frac{\mathrm{d}t}{t}
 -\frac{\mathrm{d}t}{t+1},\qquad
 \frac{\mathrm{d}y}{y-1}=-\frac{\mathrm{d}t}{t+1}.
\end{equation}
Its derivative at $t=0$ is one, so the logarithmic basepoint adds no
constant: $\ln y=\ln t-\ln(1+t)$.  Applying these replacements to the
independent angular primitive, including its rational coefficients,
gives exactly thirteen words:
\begin{equation}
 C_{V+\Sigma}(s,\lambda)
 =\sum_{w\in\mathcal W_V}v_w(X,\widehat c_Z)G(w;X),\qquad
 \widehat c_Z=-4-4\ln\lambda.
\end{equation}
Here $\lambda$ is measured in units of $m$. The symbol
$\widehat c_Z$ is the ultraviolet-finite field coefficient used in the
angular primitive, not the full counterterm $c_Z$ of Chapter 2.
With $G(\varnothing;X)=1$, the word set is
\begin{align}
 \mathcal W_V={}&\{\varnothing,(-1),(0),(-1,-1),(-1,0),(0,-1),(1,0),\\
 &(-1,-1,-1),(-1,-1,0),(-1,0,-1),(0,-1,-1),(0,-1,0),(0,1,0)\}.
\end{align}
Every $v_w$ is rational in $X$ and $\widehat c_Z$.  All explicit $\pi^2$ constants
cancel in this representation.  For example,
\begin{align}
 v_{(0,1,0)}&=v_{(0,-1,0)}=\frac{4(X+1)}{X^3},\\
 v_{(-1,-1,0)}&=-\frac{4(X-1)}{X^3},\qquad
 v_{(1,0)}=\frac{X^2-9X+16}{X^3}.
\end{align}
The full coefficient and derivative dictionaries are obtained from the
independent trace and angular-primitive files, without using the
published total.  This reduces the complete vertex and self-energy subset. It is not
separately gauge invariant and must be combined with the spanning boxes.

\subsection{The direct spanning box}
The direct spanning-box channel can also be reduced completely.  In its
angular integral let
\begin{equation}
 b=\sqrt{\frac{X}{X+4}},\qquad
 y=\frac{X-(X+4)v^2}{1-v^2},\qquad
 0<v<V=\frac{X}{X+2}.
\end{equation}
Its input GPL words use only $0,\pm1$.  Exact rational integration in
$v$ adds first letters $\pm b$, with no other roots.  Setting
$X=4b^2/(1-b^2)$ makes every rational coefficient and partial-fraction
pole rational in $b$.  The resulting angular primitive has been
checked by differentiation, with every GPL coefficient canceling
identically for arbitrary $b$.

At the physical endpoint the variable $q$ gives
\begin{equation}
 b=\frac{1-q}{1+q},\qquad V=\frac{(1-q)^2}{1+q^2}.
\end{equation}
Both quantities are rational in $q$.  Let $\mathcal G$ denote GPLs
based at $q=1$, with the same real left tangent used above.  The
weight-one transformations are
\begin{align}
 G(0;V)&=2\mathcal G(1;q)-\mathcal G(\ii;q)
                   -\mathcal G(-\ii;q)-\ln2,\\
 G(1;V)&=\mathcal G(0;q)-\mathcal G(\ii;q)-\mathcal G(-\ii;q),\\
 G(-1;V)&=\mathcal G(\rho_+;q)+\mathcal G(\rho_-;q)
                  -\mathcal G(\ii;q)-\mathcal G(-\ii;q),\\
 G(b;V)&=2\mathcal G(0;q)-\mathcal G(\ii;q)-\mathcal G(-\ii;q),\\
 G(-b;V)&=-\mathcal G(\ii;q)-\mathcal G(-\ii;q).
\end{align}
These follow by taking logarithms of the corresponding rational ratios,
with their constants fixed at $q=1$.  In particular, for
$q=1-\delta$ one has $V=\delta^2/2+O(\delta^3)$, so the $-\ln2$
in the first line is necessary.

Higher weights follow recursively from the total differential of a
moving-letter GPL.  All differences between $0,\pm1,\pm b$ and $V$
are rational functions of $q$.  Their logarithmic derivatives have
simple poles only in
\begin{equation}
 \mathcal A_s=\{0,1,-1,\ii,-\ii,\rho_+,\rho_-\}.
\end{equation}
After a lower-weight word has been converted, integration of each term
$\nu_a\mathcal G(w;q)$ prepends the letter $a$.  Every original word
containing a nonzero letter vanishes as $q\to1$, since $V=O(\delta^2)$
while $b=O(\delta)$.  Thus its integration constant is zero.  The
all-zero words are fixed by powers of the first weight-one identity.
Coincident adjacent letters cancel in the total differential before
individual logarithmic derivatives are taken.  The transformation
checks these cancellations exactly, including the cancellation of
any spurious pure powers of $\ln(1-q)$ at the boundary.

Finally, the scalar factors multiplying the angular words satisfy
\begin{equation}
 d_s=\frac{G(-1,0;X)}{X},\qquad \ell_X=G(0;X).
\end{equation}
Transforming these factors with the kernel map above and applying
shuffle identities completes the direct-box result as 101 constant-letter
words of weight at most three.  Its rational coefficients depend only
on $q$ and $\ln\lambda$, with possible constant $\ln2$ factors.  The
normalization is $C_{\mathrm{box},s}=J_s/(2X^2)$.
This reduction uses only the independent box coefficients.

\subsection{The crossed spanning box in the energy basis}
Section~\ref{sec:comp-crossed-first} has already integrated the rational
crossed subgroup in $y$, before introducing the angular square root.
Its endpoint maps contain only the letters $-1,0,1$. Applying the
kernel pullbacks above converts those words to the same $q$ basis.
The remaining angular integral has no auxiliary $\rho_u$ letter.

\subsection{The complete virtual contribution}
After the crossed subgroup is integrated as just described, its remaining
seventeen angular words involve only $0,\pm1,\pm b$.  The same total
differential construction used for the direct box therefore applies.
Combining it with the crossed subgroup gives 110 constant-letter words
for $J_u$.  The two independent constructions of the subgroup endpoint
coefficients agree as exact rational functions.  The remaining angular
primitive has also been checked by an exact derivative, coefficient by
coefficient, for arbitrary kinematics.

The entire independently computed virtual contribution is consequently
in the comparison basis:
\begin{equation}
 C_V(q,\lambda)=\sum_{w\in\mathcal W}V_w(q,\ln\lambda)\mathcal G(w;q).
\end{equation}
There are 116 nonzero words.  All explicit $\ln2$ coefficients introduced
by intermediate changes of endpoint cancel in this complete virtual
sum.  This is an algebraic cancellation of coefficients, not a numerical
inference.  The coefficients still depend on $\ln\lambda$, as expected
for the virtual contribution.

\subsection{Exact inversion of the kernel transformation}
The common basis can be reduced without choosing functions from a
published answer. Let $M$ be the $9\times6$ matrix of coefficients
in the six kernel pullbacks above. A row labels a letter in
$\mathcal A$, and a column labels one of the six $X$ kernels. Its
six columns are linearly independent. Select six independent rows
with a projection matrix $P$ and define
\begin{equation}
 L=(PM)^{-1}P,\qquad LM=I_6.
\end{equation}
For weight $n$, let $v_n$ denote the entire coefficient vector of
the independently integrated $q$ words, with zero entries for absent
words. Applying the tensor product of $n$ copies of $L$ gives the
candidate $X$ coefficients $c_n$. The necessary and sufficient
reconstruction test is
\begin{equation}
 c_n=L^{\otimes n}v_n,\qquad
 M^{\otimes n}c_n-v_n=0,\qquad 0\leq n\leq3.
\end{equation}
Every component of the second equation is checked as an exact
rational identity, including the rows not selected by $P$. Thus the
choice of independent rows cannot hide an omitted word. The
weight-zero equation includes constants. The common basepoint and
logarithmic normalization established above make this an equality
of integrals, including their boundary constants.

The remaining coefficient functions are rational in $q$. Polynomial
reduction using $q^2-(X+2)q+1=0$ reduces each one to $A(X)+qB(X)$.
The physical root is fixed by $0<q<1$:
\begin{equation}
 q=\frac{X+2-\sqrt{X(X+4)}}2.
\end{equation}
The other root exceeds one for $X>0$. Substitution therefore gives
coefficients rational in $X$ and $Y=\sqrt{X/(X+4)}$, with no branch
choice left to numerical evaluation.
